% Unnumbered introduction of the supplement.
\noindent This supplement gives the formal model behind the formalisation section of the main
text, the full proofs of time-travel stability, of the soundness proposition and of the
exactness proposition, the refinements for counting and latest-state selection, the list of
fragment assumptions, and the full grammar of the intermediate representation (IR).
Section~S1 fixes the model, S2 proves time-travel stability, S3 soundness, S4 exactness and the
over-approximation examples, S5 lists every assumption the proofs use, and S6 gives the IR
grammar and type rules.

\paragraph{Status of the statements.}
Lemmas~\ref{lem:collapse}--\ref{lem:evidence}, Propositions~\ref{prop:s-soundness}
and~\ref{prop:s-exactness}, Corollaries~\ref{cor:stepwise}--\ref{cor:filters} and
Remarks~\ref{rem:count} and~\ref{rem:latest} are proved here for the model of S1. That the
implementation realises the model (the compiled SPARQL computes the rows of S1, and the change
matcher decides the triggers of S3) is not proved: it is assumption~A3 and~A9 of S5, checked by
the test suite of the released code and, on gold graphs, by counting the misses of every sound
strategy against read-time re-query (ablation A). No statement here is a conjecture; where a
claim is only illustrated by an example, it is labelled as an example.

\paragraph{Credit.}
The compositional argument is that of Proposition~1 of Martin, Unbehauen and
Auer~\cite{martin2010caching}; the proposition is an instance of query--update
independence~\cite{blakeley1986updating,levy1993independent}. What is added is the treatment of
negation, distinct counting, latest-state selection, cut-off reads on an append-only bitemporal
graph, claim patterns with attribution constraints as the matching unit, and the witnesses and
evidence of a view, together with the exactness conditions. Each step is routine.
